\section{Gaussian Space --- Sum Invariance}
%%% generated by the blueprint pipeline from TCSlib.BooleanAnalysis.GaussianSpace.SumInvariance
\subsection*{Gaussian Space --- Sum Invariance: Introduction}

Declarations of the Lean module \texttt{TCSlib.BooleanAnalysis.GaussianSpace.SumInvariance}.

\subsection*{Gaussian Space --- Sum Invariance: Results}

\begin{definition}[Sum of third moments]
\lean{BooleanAnalysis.GaussianSpace.thirdMomentSum}
\statementsource{Analysis-of-Boolean-Functions-by-Ryan-ODonnell}{pdf-ffbe9c7b3306-p119-b002}
\label{BooleanAnalysis.GaussianSpace.thirdMomentSum}
\leanok
Let $n$ be a natural number, let $(\Omega,\Sigma)$ be a measurable space with measure
$\mu$, and let $X_1,\ldots,X_n:\Omega\to\mathbb{R}$ be measurable-space functions. The
third-moment sum is $\sum_{i=1}^n\int_\Omega |X_i(\omega)|^3\,d\mu(\omega)$.
\end{definition}

\begin{definition}[Families with matching first two moments]
\lean{BooleanAnalysis.GaussianSpace.MatchesSumMoments}
\label{BooleanAnalysis.GaussianSpace.MatchesSumMoments}
\leanok
For a natural number $n$, measurable spaces $\Omega$ and $\Xi$, measurable-space
measures $\mu$ and $\nu$, and real-valued families $(X_i)_{i<n}$ on $\Omega$
and $(Y_i)_{i<n}$ on $\Xi$, the families are each independent and every member has a
finite third moment. For each $i<n$, the integrals of $X_i$ and $Y_i$ agree, and the
integrals of $X_i^2$ and $Y_i^2$ agree; no zero-mean or normalization condition is
imposed.
\end{definition}

\begin{definition}[Lévy distance between measures]
\lean{BooleanAnalysis.GaussianSpace.levyDistance}
\label{BooleanAnalysis.GaussianSpace.levyDistance}
\leanok
For measures $\mu$ and $\nu$ on $\mathbb R$, define their Lévy distance to be the
infimum of all real numbers $\varepsilon \ge 0$ such that, for every $u \in \mathbb R$,
the measure of $(-\infty,u-\varepsilon]$ under $\nu$, converted to a real number and
then decreased by $\varepsilon$, is at most the measure of $(-\infty,u]$ under $\mu$,
converted to a real number, and that latter quantity is at most the measure of
$(-\infty,u+\varepsilon]$ under $\nu$, converted to a real number and then increased by
$\varepsilon$.
\end{definition}

\begin{theorem}[Smooth sum invariance bound]
\lean{BooleanAnalysis.GaussianSpace.sum_invariance}
\label{BooleanAnalysis.GaussianSpace.sum_invariance}
\leanok
\uses{BooleanAnalysis.GaussianSpace.MatchesSumMoments, BooleanAnalysis.GaussianSpace.thirdMomentSum}
Let $n\in\mathbb N$, let $X_1,\ldots,X_n$ and $Y_1,\ldots,Y_n$ be real-valued random
variables on (possibly different) probability spaces, with each family independent.
Assume every variable has finite third absolute moment, and for each $i$ the means and
second moments of $X_i$ and $Y_i$ agree. If $\psi\colon\mathbb R\to\mathbb R$ is three
times continuously differentiable and $C\ge0$ satisfies $|\psi'''(t)|\le C$ for every
real $t$, then the absolute difference between the expectations of $\psi(\sum_i X_i)$
and $\psi(\sum_i Y_i)$ is at most $\frac C6\sum_i\bigl(\mathbb E|X_i|^3+\mathbb
E|Y_i|^3\bigr)$.
\end{theorem}

\begin{theorem}[Berry–Esseen bound for smooth functions]
\lean{BooleanAnalysis.GaussianSpace.variant_berryEsseen}
\label{BooleanAnalysis.GaussianSpace.variant_berryEsseen}
\leanok
\uses{BooleanAnalysis.GaussianSpace.thirdMomentSum, BooleanAnalysis.ThresholdFunctions.standardGaussianMeasure}
Let $n\in\mathbb N$, let $(\Omega,\mu)$ be a probability space, and let $X_1,\ldots,X_n$
be independent real random variables on it, each in $L^3(\mu)$. Suppose each has mean
zero and their second moments sum to $1$. If $\psi:\mathbb R\to\mathbb R$ is three times
continuously differentiable, and $C\ge0$ satisfies $|\psi'''(t)|\le C$ for every
$t\in\mathbb R$, then
\[
\left|\int_\Omega\psi\!\left(\sum_{i=1}^n X_i\right)d\mu-\int_{\mathbb
R}\psi(z)\,d\gamma(z)\right|
\le \frac{1+2\sqrt{2/\pi}}{6}\,C\sum_{i=1}^n\int_\Omega|X_i|^3\,d\mu,
\]
where $\gamma$ is the standard Gaussian probability measure.
\end{theorem}

\begin{theorem}[Smooth Lipschitz approximation with derivative bounds]
\lean{BooleanAnalysis.GaussianSpace.exists_smooth_lipschitz_approximation}
\label{BooleanAnalysis.GaussianSpace.exists_smooth_lipschitz_approximation}
\leanok
There is a sequence of positive real constants $(C_k)_{k\in\mathbb N}$. For every $c\ge
0$, every $c$-Lipschitz function $\psi\colon\mathbb R\to\mathbb R$, and every $\eta>0$,
there is a function $\psi_s\colon\mathbb R\to\mathbb R$ that is infinitely
differentiable, satisfies $|\psi(t)-\psi_s(t)|\le c\eta$ for every $t\in\mathbb R$, and
satisfies $|\psi_s^{(k)}(t)|\le C_k c/\eta^{k-1}$ for every integer $k\ge 1$ and every
$t\in\mathbb R$.
\end{theorem}

\begin{theorem}[Lipschitz sum invariance bound]
\lean{BooleanAnalysis.GaussianSpace.sum_invariance_lipschitz}
\label{BooleanAnalysis.GaussianSpace.sum_invariance_lipschitz}
\leanok
\uses{BooleanAnalysis.GaussianSpace.MatchesSumMoments, BooleanAnalysis.GaussianSpace.thirdMomentSum}
There is a constant $K>0$ such that for every natural number $n$, every two probability
spaces, and independent real-valued families $X_1,\ldots,X_n$ and $Y_1,\ldots,Y_n$ on
those respective spaces, the following holds whenever each variable has finite third
absolute moment and corresponding variables have equal means and equal second moments.
For every $c\ge0$ and every function $\psi\colon\mathbb R\to\mathbb R$ that is Lipschitz
with constant at most $c$, the absolute difference between the expectations of
$\psi(\sum_iX_i)$ and $\psi(\sum_iY_i)$ is at most $Kc\bigl(\sum_i(\mathbb
E|X_i|^3+\mathbb E|Y_i|^3)\bigr)^{1/3}$.
\end{theorem}

\begin{theorem}[Lévy distance bound for sums]
\lean{BooleanAnalysis.GaussianSpace.sum_invariance_levy}
\label{BooleanAnalysis.GaussianSpace.sum_invariance_levy}
\leanok
\uses{BooleanAnalysis.GaussianSpace.MatchesSumMoments, BooleanAnalysis.GaussianSpace.levyDistance, BooleanAnalysis.GaussianSpace.thirdMomentSum}
There is a constant $K>0$ such that the following holds for every natural number $n$,
every two probability spaces, and independent real-valued families $X_1,\ldots,X_n$ and
$Y_1,\ldots,Y_n$ on those respective spaces. If each variable has finite third absolute
moment and corresponding variables have equal means and equal second moments, then the
Lévy distance between the distributions of $\sum_i X_i$ and $\sum_i Y_i$ is at most
$K\bigl(\sum_i(\mathbb E|X_i|^3+\mathbb E|Y_i|^3)\bigr)^{1/4}$. Here the Lévy distance
between probability measures $\mu,\nu$ on $\mathbb R$ is the infimum of the numbers
$\varepsilon\ge0$ such that for every $u\in\mathbb R$, $\nu(( -\infty,u-\varepsilon])
-\varepsilon\le\mu(( -\infty,u])\le\nu(( -\infty,u+\varepsilon]) +\varepsilon$.
\end{theorem}

\begin{theorem}[Berry–Esseen bound with a quarter-power rate]
\lean{BooleanAnalysis.GaussianSpace.berryEsseen_quarter_power}
\label{BooleanAnalysis.GaussianSpace.berryEsseen_quarter_power}
\leanok
\uses{BooleanAnalysis.GaussianSpace.thirdMomentSum, BooleanAnalysis.ThresholdFunctions.kolmogorovDistance, BooleanAnalysis.ThresholdFunctions.standardGaussianMeasure}
There is a constant $K>0$ such that the following holds. Let $X_1,\ldots,X_n$ be
independent real random variables on a probability space, each with finite third
absolute moment, with $\mathbb E X_i=0$ and $\sum_{i=1}^n\mathbb E[X_i^2]=1$. The
Kolmogorov distance between the law of $\sum_{i=1}^n X_i$ and the standard Gaussian
measure is at most $K\bigl(\sum_{i=1}^n\mathbb E|X_i|^3\bigr)^{1/4}$, with the same $K$
for every such family.
\end{theorem}
